\Needspace{6\baselineskip}
\Article{326}
The preceding method is already so convenient in itself that it leaves hardly anything to be desired. Nevertheless, various devices can shorten it considerably, and we can touch on only a few of them here. We shall therefore restrict the investigation to the case where the exclusion number is an odd prime not dividing \(A\), or a power of such a prime, especially since the other cases can either be reduced to this one or treated by an analogous method. Suppose, \textit{first}, that the exclusion number \(E=p\) is a prime dividing neither of the numbers \(m,\allowbreak n\), and that the values of the expressions \(\frac Am,\allowbreak -\frac{na}m,\allowbreak -\frac{nb}m\), \(-\frac{nc}m\), etc., \((\mathrm{mod.}\ p)\) are respectively \(k,\allowbreak \mathfrak A,\allowbreak \mathfrak B,\allowbreak \mathfrak C\), etc. The numbers \(\alpha,\allowbreak \beta,\allowbreak \gamma\), etc., are then found from the congruences \(\alpha\equiv k+\mathfrak A,\allowbreak \ \beta\equiv k+\mathfrak B,\allowbreak \ \gamma\equiv k+\mathfrak C\), etc., \((\mathrm{mod.}\ p)\). But the numbers \(\mathfrak A,\allowbreak \mathfrak B,\allowbreak \mathfrak C\), etc., can be obtained without calculating the congruences, by a device entirely similar to that used in Article \(322\). They will coincide either with all nonresidues or with all residues of \(p\), except \(0\), according as the value of the expression \(-\frac mn\ (\mathrm{mod.}\ p)\), or, equivalently here, the number \(-mn\), is a residue or a nonresidue of \(p\). Thus, in \textit{example} II of the preceding article, for \(E=17\) we have \(k=7\); \(-mn=-1365\equiv12\) is a nonresidue of \(17\). Hence the numbers \(\mathfrak A,\allowbreak \mathfrak B\), etc., will be \(1,\allowbreak 2,\allowbreak 4,\allowbreak 8,\allowbreak 9,\allowbreak 13,\allowbreak 15,\allowbreak 16\), and the numbers \(\alpha,\allowbreak \beta\), etc., \(8,\allowbreak 9,\allowbreak 11,\allowbreak 15,\allowbreak 16,\allowbreak 3,\allowbreak 5,\allowbreak 6\). Of these, the residues are \(8,\allowbreak 9,\allowbreak 15,\allowbreak 16\), giving \(\pm h,\allowbreak h'\), etc., the values \(\pm5,\allowbreak 3,\allowbreak 7,\allowbreak 4\). — Those who often have occasion to solve such problems will make their work much easier if, for several primes \(p\), they calculate the values of \(h,\allowbreak h'\), etc., corresponding to the individual values of \(k\ (1,\allowbreak 2,\allowbreak 3\ldots p-1)\), under each of the two assumptions that \(-mn\) is a residue or a nonresidue of \(p\). We also observe that the number of values \(h,\allowbreak -h,\allowbreak h'\), etc., is always \(\frac12(p-1)\) when both \(k\) and \(-mn\) are residues or both are nonresidues of \(p\); \(\frac12(p-3)\) when the former is \(R.\) and the latter \(NR.\); and \(\frac12(p+1)\) when the former is \(NR.\) and the latter \(R.\). But we must omit the proof of this theorem to avoid becoming too lengthy.

\textit{Second}, the cases where \(E\) is a prime dividing \(n\), or a power of an odd prime that does or does not divide \(n\), can be treated still more conveniently. We shall include all these cases together and, retaining all the notation of Article \(324\), put \(n=n'p^\nu\), so that \(n'\) is not divisible by \(p\). The numbers \(a,\allowbreak b,\allowbreak c\), etc., will be the products of \(p^{\mu-1}\) either by all numbers less than \(p\), except \(0\), or by all nonresidues of \(p\) below \(p\), according as \(\mu\) is even or odd; denote them generally by \(up^{\mu-1}\). Let \(k\) be the value of the expression \(\frac Am\ (\mathrm{mod.}\ p^{\mu+\nu})\). It will not be divisible by \(p\), since the same property is assumed of \(A\);
\SourceLeaf{402}
further, it is clear that all of \(\alpha,\allowbreak \beta,\allowbreak \gamma\), etc., are congruent to \(k\) modulo \(p\), and hence that \(p^\mu\) excludes nothing from \(\Omega\) if \(kNp\). If, however, \(kRp\), and hence also \(kRp^{\mu+\nu}\), let \(r\) be a value of the expression \(\surd k\ (\mathrm{mod.}\ p^{\mu+\nu})\), which will not be divisible by \(p\), and \(e\) the value of \(-\frac{n'}{2mr}\ (\mathrm{mod.}\ p)\). Then \(\alpha\equiv rr+2era p^\nu\ (\mathrm{mod.}\ p^{\mu+\nu})\), from which it follows easily that \(\alpha\) is a residue of \(p^{\mu+\nu}\), and that the values of the expression \(\surd\alpha\ (\mathrm{mod.}\ p^{\mu+\nu})\) are \(\pm(r+ea p^\nu)\); hence all of \(h,\allowbreak h',\allowbreak h''\), etc., are expressed by \(r+uep^{\mu+\nu-1}\). Finally, it follows without difficulty that the numbers \(h,\allowbreak h',\allowbreak h''\), etc., arise by adding \(r\) to the products of \(p^{\mu+\nu-1}\) either by all numbers below \(p\), except \(0\), when \(\mu\) is even; or by all nonresidues of \(p\) below this bound, when \(\mu\) is odd and \(eRp\), or, equivalently here, when \(-2mrn'Rp\); or by all residues, except \(0\), when \(\mu\) is odd and \(-2mrn'Np\).

\begin{samepage}
Moreover, once the numbers \(h,\allowbreak h'\), etc., have been found for each exclusion number we wish to use, the exclusion itself can also be performed by mechanical operations, which anyone familiar with these matters can readily devise for themselves if it seems worth the effort.
\par\end{samepage}

Finally, we must observe that any equation \(axx+2bxy+cyy=M\) in which \(bb-ac\) is negative, \(=-D\), can easily be reduced to the form considered above. For, denoting the greatest common divisor of \(a,\allowbreak b\) by \(m\), and putting \[ a=ma'.\quad b=mb',\quad \frac Dm=a'c-mb'b'=n,\quad a'x+b'y=x' \] the equation is evidently equivalent to \(mx'x'+nyy=a'M\), which can be solved by the rules given above. Of its solutions, however, we must retain only those in which \(x'-b'y\) is divisible by \(a'\), or from which \(x\) acquires integer values.

\Needspace{9\baselineskip}\begin{center}\textit{Another method for solving the congruence \(xx\equiv A\) when \(A\) is negative.}\end{center}


\Needspace{6\baselineskip}\Article{327}
Just as the direct solution of the equation \(axx+2bxy+cyy=M\) given in Section V assumes the values of the expression \(\surd{(bb-ac)}\ (\mathrm{mod.}\ M)\) to be known, so, conversely, when \(bb-ac\) is negative, the indirect solution explained above supplies a very convenient method for finding these values. Especially for a very large value of \(M\), it is far preferable to the method of Article \(322\) and the following articles. We shall assume, however, that \(M\) is a prime number, or at least that its factors, if it is composite,
\SourceLeaf{403}
are still unknown. For, if we knew that a prime \(p\) divides \(M\), and that \(M=p^\mu M'\), with \(M'\) no longer containing the factor \(p\), it would be much more convenient to find the values of the expression \(\surd{(bb-ac)}\) separately for the moduli \(p^\mu\) and \(M'\), the former from the values modulo \(p\), Article \(101\), and to deduce the values modulo \(M\) by combining these, Article \(105\).

We are therefore to find all the values of the expression \(\surd{-D}\ (\mathrm{mod.}\ M)\), assuming \(D\) and \(M\) positive, and \(M\) of a form of divisors of \(xx+D\), Article \(147\) and the following articles; otherwise it would already be certain that no numbers could satisfy the proposed expression. Let the required values, which always occur in opposite pairs, be \(\pm r,\allowbreak \pm r',\allowbreak \pm r''\), etc., and let \(D+rr=Mh,\allowbreak \ D+r'r'=Mh',\allowbreak \ D+r''r''=Mh''\), etc. Further, denote the classes to which the forms \((M,\allowbreak r,\allowbreak h),\allowbreak \ (M,\allowbreak -r,\allowbreak h),\allowbreak \ (M,\allowbreak r',\allowbreak h')\), \((M,\allowbreak -r',\allowbreak h'),\allowbreak \ (M,\allowbreak r'',\allowbreak h'').\ (M,\allowbreak -r'',\allowbreak h'')\), etc., belong respectively by \(\mathfrak C,\allowbreak -\mathfrak C,\allowbreak \mathfrak C',\allowbreak -\mathfrak C'\), \(\mathfrak C'',\allowbreak -\mathfrak C''\), etc., and their collection by \(\mathfrak G\). Generally speaking, these classes must be regarded as unknown. Nevertheless, it is clear, \textit{first}, that all are positive and properly primitive; \textit{second}, that all belong to the same genus, whose \textit{character} can easily be determined from the nature of the number \(M\), \textit{i.e.} from its relations to the individual prime divisors of \(D\), and also to \(4\) or \(8\) when these are necessary, Article \(230\). Since \(M\) is assumed to be of a form of divisors of \(xx+D\), we can be certain beforehand that this character must correspond to a genus of positive, properly primitive forms of determinant \(-D\), even if the expression \(\surd{-D}\ (\mathrm{mod.}\ M)\) should perhaps have no solutions. Thus, since this genus is known, all the classes it contains can be found; let them be \(C,\allowbreak C',\allowbreak C''\), etc., and their collection \(G\).

It is therefore clear that each class \(\mathfrak C,\allowbreak -\mathfrak C\), etc., must be identical with some class in \(G\). Several classes in \(\mathfrak G\) may also be identical with one another, and hence with the same class in \(G\); and when \(G\) contains just one class, all those in \(\mathfrak G\) certainly coincide with it. Thus, if, from the classes \(C,\allowbreak C',\allowbreak C''\), etc., we select the simplest forms \(f,\allowbreak f',\allowbreak f''\), etc., one from each, a form from each class in \(\mathfrak G\) will be found among them. Now, if \(axx+2bxy+cyy\) is a form contained in the class \(\mathfrak C\), there will be two representations of \(M\) by it belonging to the value \(r\); if one is \(x=m,\allowbreak y=n\), the other is \(x=-m,\allowbreak y=-n\). The sole exception is the case \(D=1\), in which there are four representations; see Article \(180\).

It follows that, if all representations of \(M\) by each of the forms \(f,\allowbreak f',\allowbreak f''\), etc., are found by the indirect method given above, and
\SourceLeaf{404}
the values of the expression \(\surd{-D}\ (\mathrm{mod.}\ M)\) to which the individual representations belong are then deduced, Article 154 and the following articles, \textit{all} values of this expression will be obtained, each twice, or, if \(D=1\), four times. \textit{Q. E. F.} If there are forms among \(f,\allowbreak f'\), etc., by which \(M\) cannot be represented, this indicates that they belong to no class in \(\mathfrak G\), and they must therefore be disregarded. If, however, \(M\) cannot be represented by any of these forms, \(-D\) must be a quadratic nonresidue of \(M\). --- The following observations should also be kept in mind in these operations.

I. The representations of \(M\) by the forms \(f,\allowbreak f'\), etc., used here are understood to be those in which the values of the indeterminates are coprime. If others occur in which these values have a common divisor \(\mu\), which can happen only when \(\mu\mu\) divides \(M\) and will certainly happen when \(-D\,\mathrm R\,\frac{M}{\mu\mu}\), they must be entirely disregarded for the present purpose, though they may be useful in another respect.

II. Other things being equal, the work is evidently easier the smaller the number of classes \(f,\allowbreak f',\allowbreak f''\), etc., and therefore shortest when \(D\) is one of the 65 numbers given in Article 303, for which each genus contains only one class.

III. Since two such representations \(x=m,\allowbreak y=n\) and \(x=-m,\allowbreak y=-n\) always belong to the same value, it is clear that it is enough to consider only representations in which \(y\) is positive. Distinct representations \textit{of this kind} therefore always correspond to distinct values of the expression \(\surd{-D}\ (\mathrm{mod.}\ M)\), so that the number of all distinct values equals the number of all resulting representations of this kind, always excepting \(D=1\), when the former number is half the latter.

IV. Since knowing either of the two opposite values \(+r,\allowbreak -r\) immediately gives the other, the operations can be shortened somewhat further. If the value \(r\) is obtained from a representation of \(M\) by a form contained in the class \(C\), \textit{i.e.} if \(\mathfrak C=C\), the opposite value \(-r\) will evidently arise from a representation by a form in the class opposite to \(C\), which differs from the class \(C\) unless the latter is ambiguous. Hence, when not all classes in \(G\) are ambiguous, only half the remaining classes need be considered: one from each opposite pair. The other can be disregarded, since, without calculation, we can already foresee that its values will be opposite to those supplied by the first. When \(C\) is ambiguous, however, both values \(r\) and \(-r\) arise from it together. Namely, if the ambiguous form selected from \(C\) is
\SourceLeaf{405}
\(axx+2bxy+cyy\), and the value \(r\) came from the representation \(x=m,\allowbreak y=n\), the value \(-r\) will come from \(x=-m-\frac{2bn}{a},\allowbreak \ y=n\).

\begin{samepage}
V. In the case \(D=1\), there is just one class altogether, from which we may suppose that the form \(xx+yy\) has been selected. If the value \(r\) comes from the representation \(x=m,\allowbreak y=n\), the same value comes from \(x=-m,\allowbreak y=-n;\allowbreak \ x=n,\allowbreak y=-m;\allowbreak \ x=-n,\allowbreak y=m\), and the opposite value \(-r\) from \(x=m,\allowbreak y=-n;\allowbreak x=-m,\allowbreak y=n;\allowbreak \ x=n,\allowbreak y=m;\allowbreak \ x=-n,\allowbreak y=-m\). Thus one of these eight representations, which constitute a single decomposition, is enough, provided we associate the opposite value with the one obtained.
\par\end{samepage}

VI. By Article 155, the value of the expression \(\surd{-D}\ (\mathrm{mod.}\ M)\) to which the representation \(M=amm+2bmn+cnn\) belongs is \(\mu(mb+nc)-\nu(ma+nb)\), or any number congruent to it modulo \(M\), with \(\mu,\allowbreak \nu\) chosen so that \(\mu m+\nu n=1\). Denoting such a value by \(v\), we have \[\begin{aligned}
mv&\equiv\mu m(mb+nc)-\nu(M-mnb-nnc)\\
&\equiv(\mu m+\nu n)(mb+nc)\equiv mb+nc\ (\mathrm{mod.}\ M)
\end{aligned}\] Thus \(v\) is a value of the expression \(\frac{mb+nc}{m}\ (\mathrm{mod.}\ M)\); similarly, we find that \(v\) is a value of the expression \(-\frac{ma+nb}{n}\ (\mathrm{mod.}\ M)\). These formulas are often preferable to the one from which they were deduced.

\Needspace{6\baselineskip}\Article{328}
\textit{Examples.} I. We seek all values of the expression \(\surd{-1365}\ (\mathrm{mod.}\ 5428681=M)\). The number \(M\) is here \(\equiv1,\allowbreak 1,\allowbreak 1,\allowbreak 6,\allowbreak 11\ (\mathrm{mod.}\ 4,\allowbreak 3,\allowbreak 5,\allowbreak 7,\allowbreak 13)\), and is therefore of a form of divisors of \(xx+1,\allowbreak \ xx+3,\allowbreak \ xx-5\), and of a form of non-divisors of \(xx+7,\allowbreak \ xx-13\), and hence of a form of divisors of \(xx+1365\). The character of the genus in which the classes \(\mathfrak G\) are found will be \(1,\allowbreak 4;\allowbreak \ R3;\allowbreak \ R5;\allowbreak \ N7;\allowbreak \ N13\). This genus contains just one class, from which we choose the form \(6xx+6xy+229yy\). To find all representations of \(M\) by it, we put \(2x+y=x'\), giving \(3x'x'+455yy=2M\). This equation admits four solutions in which \(y\) is positive, namely \(y=127,\allowbreak \ x'=\pm1083,\allowbreak \ y=119,\allowbreak \ x'=\pm1213\). These give four solutions of the equation \(6xx+6xy+229yy=M\) in which \(y\) is positive: \[\renewcommand{\arraystretch}{1.25}\begin{array}{c|r|r|r|r}x&478&-605&547&-666\\y&127&127&119&119\end{array}\] The first solution gives for \(v\) the value of the expression \(\frac{30517}{478}\), or \(-\frac{3249}{127}\ (\mathrm{mod.}\ M)\), from which we find
\SourceLeaf{406}
\(2350978\); the second produces the opposite value \(-2350978\); the third gives \(2600262\), and the fourth its opposite \(-2600262\).

II. If the values of the expression \(\surd{-286}\ (\mathrm{mod.}\ 4272943=M)\) are sought, the character of the genus containing the classes \(\mathfrak G\) is found to be \(1\text{ and }7,\allowbreak 8;\allowbreak \ R11;\allowbreak \ R13\). It is therefore the principal genus, containing three classes represented by the forms \((1,\allowbreak 0,\allowbreak 286)\), \((14,\allowbreak 6,\allowbreak 23),\allowbreak \ (14,\allowbreak -6,\allowbreak 23)\). The third may be disregarded, since it is opposite to the second. The form \(xx+286yy\) gives two representations of \(M\) in which \(y\) is positive, namely \(y=103,\allowbreak \ x=\pm1113\), yielding the values \(1493445,\allowbreak \ -1493445\) of the proposed expression. But by the form \((14,\allowbreak 6,\allowbreak 23)\), \(M\) is found not to be representable, so there are no values other than the two found.

\begin{samepage}
III. Given the expression \(\surd{-70}\ (\mathrm{mod.}\ 997331)\), the classes \(\mathfrak G\) must lie in the genus whose character is \(3\text{ and }5,\allowbreak 8;\allowbreak \ R5;\allowbreak \ N7\). There is just one class in this genus, represented by the form \((5,\allowbreak 0,\allowbreak 14)\). But calculation shows that \(997331\) cannot be represented by the form \((5,\allowbreak 0,\allowbreak 14)\); hence \(-70\) must be a quadratic nonresidue of that number.
\par\end{samepage}
